\documentclass[10pt,twocolumn]{article}
\usepackage[T1]{fontenc}
\usepackage[utf8]{inputenc}
\usepackage{lmodern}
\usepackage[margin=1.8cm,top=2.4cm,bottom=2cm,columnsep=0.6cm]{geometry}
\usepackage{fancyhdr}
\usepackage{titlesec}
\usepackage{booktabs}
\usepackage{amsmath,amssymb,amsfonts}
\usepackage{microtype}
\usepackage{xcolor}
\usepackage{mdframed}
\usepackage{parskip}
\usepackage{enumitem}
\usepackage{hyperref}

\definecolor{rulered}{RGB}{180,30,30}
\definecolor{boxgray}{RGB}{245,245,245}
\definecolor{keywordgray}{RGB}{100,100,100}

\pagestyle{fancy}
\fancyhf{}
\fancyhead[L]{\textsc{\small Claude Mag}}
\fancyhead[C]{\small\textit{Machine Learning Research Digest}}
\fancyhead[R]{\small 2026-11-23}
\fancyfoot[C]{\small\thepage}
\renewcommand{\headrulewidth}{0.8pt}

\titleformat{\section}{\normalsize\bfseries\color{rulered}}{}{0pt}{}%
  [\vspace{-4pt}\color{rulered}\rule{\columnwidth}{0.4pt}]
\titlespacing{\section}{0pt}{8pt}{4pt}

\hypersetup{colorlinks=true,urlcolor=rulered,linkcolor=black}

\begin{document}
\setlength{\parindent}{0pt}
\setlength{\parskip}{4pt}

{\small\color{keywordgray}\textbf{Keywords:}
  Robotics \textbullet{} World Models \textbullet{} JEPA \textbullet{} Latent Prediction \textbullet{} Robot Policies}\par\vspace{4pt}

{\LARGE\bfseries\raggedright La-JEPA: Compact Robot Policies via Latent World Modeling}\par\vspace{4pt}
{\large\itshape\raggedright A JEPA latent world model outperforms a 10$\times$-larger rival on LIBERO-Plus, with no task demonstrations for pretraining}\par\vspace{6pt}
{\small\textbf{By} La-JEPA Team (arXiv:2611.08156)}\par\vspace{2pt}
{\small\color{keywordgray}arXiv:2611.08156 \textbullet{} cs.RO, cs.LG \textbullet{} Submitted 8 Nov 2026}
\par\vspace{2pt}\noindent{\color{rulered}\rule{\columnwidth}{0.6pt}}\vspace{6pt}

La-JEPA (Latent-Action JEPA) trains a compact world model on raw robot observations using the JEPA objective---predicting future latent embeddings without decoding to pixels---then conditions lightweight manipulation policies on the world model's forward predictions.
La-JEPA-small (0.4B parameters) outperforms LeapBot-WA (10.4B) on LIBERO-Plus and achieves the highest average success on a 50-task RoboTwin 2.0 benchmark.

\begin{mdframed}[backgroundcolor=boxgray,linecolor=rulered,linewidth=1pt,
    innerleftmargin=6pt,innerrightmargin=6pt,
    innertopmargin=6pt,innerbottommargin=6pt]
\textbf{\small AT A GLANCE}\par\vspace{4pt}
\begin{description}[leftmargin=0pt,labelsep=4pt,itemsep=2pt,parsep=0pt,
    font=\small\bfseries]
  \item[Problem:] Robot manipulation policies require expensive per-task demonstrations; large parameters do not substitute for good world models.
  \item[Why it matters:] Compact, parameter-efficient models that generalize broadly reduce data and compute requirements for deploying robots.
  \item[Method:] JEPA latent world model on raw pixels; latent future-state conditioning at policy inference; action tokenizer for embodiment transfer.
  \item[Result:] La-JEPA-small (0.4B) beats LeapBot-WA (10.4B); best on LIBERO-Plus (10,030 tasks) and RoboTwin 2.0 (50 tasks).
  \item[Evidence:] Three model sizes (0.4B, 0.9B, 2.1B); two large-scale benchmarks.
\end{description}
\end{mdframed}

\section{The Big Picture}

Robot manipulation generalization has a parameter problem: the best-performing systems tend to be very large models trained on massive demonstration datasets.
La-JEPA challenges both scaling axes simultaneously: it achieves state-of-the-art results at a fraction of the parameter count of leading competitors, and it requires no task-labeled demonstrations for world-model pretraining.

The key idea is a clean separation between two learning problems.
\textit{World model pretraining} asks: what are the dynamics of the visual scene given an action?
This can be answered from raw observations with an unsupervised predictive objective.
\textit{Policy learning} asks: given a goal, what actions should I take?
This requires task-labeled demonstrations, but only a small number once the world model already understands scene dynamics.

La-JEPA uses JEPA (Joint Embedding Predictive Architecture) for the first problem and a lightweight behavioral cloning head for the second.

\section{Why Existing Approaches Fall Short}

\textbf{Large-scale VLA models} (e.g.\ RT-2, $\pi_0$) achieve generalization through scale, requiring billions of parameters and massive demonstration corpora.
Deployment is expensive; per-embodiment re-training is impractical.

\textbf{Prior JEPA-based robots} (e.g.\ VLA-JEPA, LeWorldModel) use JEPA pretraining but couple it to task-specific prompts or require a separate large pretrained encoder, adding parameter cost and limiting generalization.

\textbf{Pixel-reconstructing world models} (e.g.\ DreamerV3 adapted for manipulation) spend capacity modeling visually irrelevant background variation, leaving less capacity for task-relevant dynamics.

La-JEPA avoids pixel reconstruction entirely: the JEPA objective forces the latent space to capture task-relevant dynamics directly.

\section{Core Mechanism}

La-JEPA's world model predicts the latent embedding of a future observation given the current observation and a compact action encoding, without ever decoding back to pixels:

\begin{enumerate}[leftmargin=1.5em,itemsep=2pt]
  \item \textbf{Encode current state:} $z_t = \phi(o_t)$ maps raw observation to latent vector.
  \item \textbf{Tokenize action:} $\tilde{a} = g(a_{t:t+k})$ maps joint-space trajectory to compact action tokens.
  \item \textbf{Predict future latent:} $\hat{z}_{t+k} = f_\theta(z_t, \tilde{a})$ predicts the latent of the future state.
  \item \textbf{Condition policy:} $a_t \sim \pi_\psi(a_t \mid z_t, \hat{z}_{t+H})$ uses the world model's $H$-step prediction as a look-ahead.
\end{enumerate}

\section{Technical Formulation}

\subsection*{JEPA World Model Objective}

Training minimizes the prediction error in latent space, with a stop-gradient on the target to prevent representational collapse:
\begin{equation}
\mathcal{L}_{\mathrm{JEPA}} = \left\|\hat{z}_{t+k} - \mathrm{sg}\!\left(\phi(o_{t+k})\right)\right\|^2
\label{eq:jepa}
\end{equation}

No pixel reconstruction loss is involved.
The encoder $\phi$ learns to represent features that are predictable from past observations and actions, automatically discarding visually irrelevant variation.

\subsection*{Collapse Prevention}

The prediction error covariance across latent dimensions is monitored:
\begin{equation}
\Sigma_z = \mathbb{E}\!\left[(\hat{z} - z)(\hat{z} - z)^\top\right]
\label{eq:cov}
\end{equation}

When $\mathrm{rank}(\Sigma_z)$ falls below a threshold, a contrastive penalty is added to the loss to expand the active latent subspace, preventing the encoder from collapsing to low-dimensional representations.

\subsection*{Full Training Objective}

\begin{equation}
\mathcal{L}(\theta, \psi) = \mathcal{L}_{\mathrm{JEPA}}(\theta) + \lambda_{\mathrm{BC}}\,\mathcal{L}_{\mathrm{BC}}(\psi) + \lambda_{\mathrm{WM}}\,\mathcal{L}_{\mathrm{WM}}(\theta, \psi)
\label{eq:full}
\end{equation}

$\mathcal{L}_{\mathrm{BC}}$ is behavioral cloning on task demonstrations; $\mathcal{L}_{\mathrm{WM}}$ trains the policy to match rollouts from the world model's forward predictions.
Pretraining minimizes $\mathcal{L}_{\mathrm{JEPA}}$ alone; fine-tuning adds the downstream terms.

\section{Learning or Inference Procedure}

\textbf{Phase 1 (World model pretraining):}
Train on raw robot observations with no task labels.
Optimize $\mathcal{L}_{\mathrm{JEPA}}$ alone.
The encoder learns visual dynamics of the manipulation workspace.

\textbf{Phase 2 (Policy fine-tuning):}
Add task demonstrations.
Train $\mathcal{L}_{\mathrm{BC}} + \mathcal{L}_{\mathrm{WM}}$ with the world model optionally frozen.
The policy learns to plan in the world model's latent space.

\textbf{Inference:}
At each step, roll the world model forward $H$ steps from the current observation and planned action sequence.
Execute the action with highest predicted task-relevant latent value.

\section{What the Guarantee Says}

No formal convergence guarantee is provided.
The empirical claim is that JEPA-based latent world models learn representations sufficient to generalize across a 10,000-task manipulation distribution, and that latent-space look-ahead substantially improves sample efficiency during policy fine-tuning.

The action tokenizer is a key design choice: by encoding trajectories at the embodiment-agnostic level, the world model can be reused across robots without re-pretraining.

\section{Experimental Findings}

\begin{center}
\small
\begin{tabular}{lrr}
\toprule
\textbf{Method} & \textbf{Params} & \textbf{LIBERO-Plus} \\
\midrule
LeapBot-WA & 10.4B & 81.2\% \\
VLA-JEPA & 3.0B & 78.9\% \\
LeWorldModel & 1.8B & 76.4\% \\
La-JEPA (small) & 0.4B & 83.7\% \\
La-JEPA (base) & 0.9B & 86.1\% \\
La-JEPA (large) & 2.1B & \textbf{88.4\%} \\
\bottomrule
\end{tabular}
\end{center}

La-JEPA-small (0.4B) outperforms the 10.4B LeapBot-WA at 25$\times$ fewer parameters.
On RoboTwin 2.0, La-JEPA achieves the highest average success on both the 20-task comparison subset and the full 50-task benchmark.

\section{Ablations and Interpretation}

\begin{itemize}[leftmargin=1em,itemsep=2pt]
  \item \textbf{Without latent look-ahead:} Removing the $H$-step prediction at policy inference reduces LIBERO-Plus success by 4.3\% at the base size, confirming that world-model conditioning is load-bearing.
  \item \textbf{Pixel reconstruction vs.\ JEPA:} A matched-parameter model with a pixel-reconstruction world model achieves 79.1\% on LIBERO-Plus, 7 points below La-JEPA base, confirming that JEPA's latent compression is advantageous.
  \item \textbf{Action tokenizer:} Without the action tokenizer (direct joint-space conditioning), world-model transfer to a second robot embodiment drops 11.4\%, confirming the tokenizer's role in embodiment abstraction.
  \item \textbf{Collapse monitoring:} Disabling the rank-based contrastive penalty causes latent collapse at training step $\approx 10{,}000$; performance plateaus at 73.2\%.
\end{itemize}

\section*{Reference}
(Author list at arXiv.)
\textbf{La-JEPA: Compact Robot Policies with Latent World Modeling.}
arXiv:2611.08156. cs.RO, cs.LG. November 2026.
\url{https://arxiv.org/abs/2611.08156}

\end{document}
